\section{Introduction and Research Questions}
\label{sec:introduction}

Regular global \texttt{new} was compared with eight versions of a custom
allocator.
The custom allocator requested one large region of memory from the operating
system and then handed out pieces from that region.
Each workload was built twice from the same source, with one
one using the built-in \texttt{new} and \texttt{delete}, and 
the other redirecting every allocation in the program into the custom memory pool.

The four workload pairs profiled were matrix array, uniform matrix nodes,
nested matrix stress, and linked list.
Each pair was run at three lengths, from a few minutes to more than six hours.
Runs were made on several machines and operating systems using matching
arguments, seeds, and build settings, with the results kept separate.
Elapsed time was recorded for every run.
Peak memory, page faults, and context switches were recorded where the
profiling tools reported them.
No single workload or machine was treated as evidence that one allocator is
always faster than another.

The investigation and analysis of this project addressed four main questions:

\begin{itemize}
  \item How did each version track free memory, select a free block, split and reuse
    blocks, recombine neighboring blocks, align allocations, and account for
    metadata?
  \item What advantages and costs did those choices carry for fragmentation,
    allocation search cost, and lock contention?
  \item What did specializing a version for matrix arithmetic trade away on
    other allocation patterns?
  \item Could a simpler multi-threaded implementation be made competitive by
    pairing it with the custom allocator?
\end{itemize}

